
%%%%  Diagram macros  %%%%

\newcommand{\tp}[2]{\begin{tikzpicture}[scale=#1] #2 \end{tikzpicture}} 
\newcommand{\flip}[1]{\reflectbox{\rotatebox[]{180}{#1}}} 


\newcommand{\sigone}[1]{\tp{#1}{
\draw (1,0) -- (2,-1); 
\draw (2,0)--(1,-1);
\draw (3,0)--(3,-1);
\draw (4,0)--(4,-1);
}}

\newcommand{\sigtwo}[1]{\tp{#1}{
\draw (1,0) -- (1,-1); 
\draw (2,0)--(3,-1);
\draw (3,0)--(2,-1);
\draw (4,0)--(4,-1);
}}

\newcommand{\onecup}[8]{ %% NOT WORKING!  draw a onecup element u_{#2#3} with n=#1 %%
\begin{tikzpicture}[scale=0.368]
     \draw (0,0) -- (#1,0) ;        %% the horizontal line 
     \foreach \x in {#5,...,#6}
          \draw (\x,0) -- (\x,-1);  %% the vert lines between cup
     \foreach \x % [evaluate \x] 
          in {1,2,...,#4}
     \draw (\x,0) -- (\x,-1);       %% the vert lines before cup
     \foreach \x in {#7,...,#8}
        \draw (\x,0) -- (\x,-1);    %% the vert lines after cup
%     \draw (2.2385,0) -- (2.2385,-1);
     \draw (#2,0) .. controls (#2,-.5) and (#3,-.5) .. (#3,0) ;
 \end{tikzpicture}
 }
 
\newcommand{\onecupp}[8]{ %% NOT WORKING!  draw a onecup element u_{#2#3} with n=#1 %%
\begin{tikzpicture}[scale=0.368]
     \draw (0,0) -- (#1,0) ;        %% the horizontal line 
     \foreach \x in {#4,#5,#6,#7,#8}
          \draw (\x,0) -- (\x,-1);  %% the vert lines 
%     \foreach \x % [evaluate \x] 
%          in {1,2,...,#4}
%     \draw (\x,0) -- (\x,-1);       %% the vert lines before cup
%     \foreach \x in {#7,...,#8}
%        \draw (\x,0) -- (\x,-1);    %% the vert lines after cup
%%%     \draw (2.2385,0) -- (2.2385,-1);
     \draw (#2,0) .. controls (#2,-.5) and (#3,-.5) .. (#3,0) ;
 \end{tikzpicture}
 }

 
\newcommand{\onecupppppo}[6]{ %% draw a onecup element u_{#2#3} with n=#1=5; no symmetriser %%
\begin{tikzpicture}[scale=0.268]
     \draw (0,0) -- (#1,0) ;        %% the horizontal line 
%     \draw (#4,0) -- (#4,-1) -- (2,-2);  %% the vert lines 
     \draw (#4,0) .. controls (#4,-1) .. (2,-2);  %% the vert lines 
%     \draw (#5,0) -- (#5,-1) -- (3,-2);  %% the vert lines 
     \draw (#5,0) .. controls (#5,-1) .. (3,-2);  %% the vert lines 
%     \draw (#6,0) -- (#6,-1) -- (4,-2);  %% the vert lines 
     \draw (#6,0) .. controls (#6,-1) .. (4,-2);  %% the vert lines  
     \draw (#2,0) .. controls (#2,-.5) and (#3,-.5) .. (#3,0) ;    %% the cup
 \end{tikzpicture}
 }


 
\newcommand{\onecuppppp}[6]{ %% draw a onecup element u_{#2#3} with n=#1=5; and symmetriser %%
\begin{tikzpicture}[scale=0.268]
     \draw (0,0) -- (#1,0) ;        %% the horizontal line 
%     \draw (#4,0) -- (#4,-1) -- (2,-2);  %% the vert lines 
     \draw (#4,0) .. controls (#4,-1) .. (2,-2);  %% the vert lines 
%     \draw (#5,0) -- (#5,-1) -- (3,-2);  %% the vert lines 
     \draw (#5,0) .. controls (#5,-1) .. (3,-2);  %% the vert lines 
%     \draw (#6,0) -- (#6,-1) -- (4,-2);  %% the vert lines 
     \draw (#6,0) .. controls (#6,-1) .. (4,-2);  %% the vert lines 
%     \draw (#7,0) -- (#7,-1) -- (5,-2);  %% the vert lines 
%     \draw (#8,0) -- (#8,-1) -- (6,-2);  %% the vert lines 
%     \draw (#7,0) .. controls (#7,-1) .. (5,-2);  %% the vert lines 
%     \draw (#8,0) .. controls (#8,-1) .. (6,-2);  %% the vert lines 
\draw (2,-2) -- (2,-2.5) ; %%symmetriser leg
\draw (3,-2) -- (3,-2.5) ; %%symmetriser leg
\draw (4,-2) -- (4,-2.5) ; %%symmetriser leg
\draw (1.79,-2.5) -- (4.1,-2.5) -- (4.1,-2.92) -- (1.79,-2.92) -- (1.79,-2.5);   %% the symmetrizer 
     \draw (#2,0) .. controls (#2,-.5) and (#3,-.5) .. (#3,0) ;    %% the cup
 \end{tikzpicture}
 }

 
\newcommand{\onecupppppx}[6]{ %% draw a onecup element u_{#2#3} with n=#1  and x in Specht
\begin{tikzpicture}[scale=0.268]
     \draw (0,0) -- (#1,0) ;        %% the horizontal line 
%     \draw (#4,0) -- (#4,-1) -- (2,-2);  %% the vert lines 
     \draw (#4,0) .. controls (#4,-1) .. (2,-2);  %% the vert lines 
%     \draw (#5,0) -- (#5,-1) -- (3,-2);  %% the vert lines 
     \draw (#5,0) .. controls (#5,-1) .. (3,-2);  %% the vert lines 
%     \draw (#6,0) -- (#6,-1) -- (4,-2);  %% the vert lines 
     \draw (#6,0) .. controls (#6,-1) .. (4,-2);  %% the vert lines 
%     \draw (#7,0) -- (#7,-1) -- (5,-2);  %% the vert lines 
%     \draw (#8,0) -- (#8,-1) -- (6,-2);  %% the vert lines 
%     \draw (#7,0) .. controls (#7,-1) .. (5,-2);  %% the vert lines 
%     \draw (#8,0) .. controls (#8,-1) .. (6,-2);  %% the vert lines 
\draw (2,-2) -- (2,-2.5) ; %%symmetriser leg
\draw (3,-2) -- (4,-2.5) ; %%symmetriser leg
\draw (4,-2) -- (3,-2.5) ; %%symmetriser leg
\draw (1.79,-2.5) -- (4.1,-2.5) -- (4.1,-2.92) -- (1.79,-2.92) -- (1.79,-2.5);   %% the symmetrizer 
     \draw (#2,0) .. controls (#2,-.5) and (#3,-.5) .. (#3,0) ;    %% the cup
 \end{tikzpicture}
 }
 
 
\newcommand{\onecuppp}[8]{ %% draw a onecup element u_{#2#3} with n+1=#1 %%
\begin{tikzpicture}[scale=0.3068]
     \draw (0,0) -- (#1,0) ;        %% the horizontal line 
%     \draw (#4,0) -- (#4,-1) -- (2,-2);  %% the vert lines 
     \draw (#4,0) .. controls (#4,-1) .. (2,-2);  %% the vert lines 
%     \draw (#5,0) -- (#5,-1) -- (3,-2);  %% the vert lines 
     \draw (#5,0) .. controls (#5,-1) .. (3,-2);  %% the vert lines 
%     \draw (#6,0) -- (#6,-1) -- (4,-2);  %% the vert lines 
     \draw (#6,0) .. controls (#6,-1) .. (4,-2);  %% the vert lines 
%     \draw (#7,0) -- (#7,-1) -- (5,-2);  %% the vert lines 
%     \draw (#8,0) -- (#8,-1) -- (6,-2);  %% the vert lines 
     \draw (#7,0) .. controls (#7,-1) .. (5,-2);  %% the vert lines 
     \draw (#8,0) .. controls (#8,-1) .. (6,-2);  %% the vert lines 
\draw (2,-2) -- (2,-2.5) ; %%symmetriser leg
\draw (3,-2) -- (3,-2.5) ; %%symmetriser leg
\draw (4,-2) -- (4,-2.5) ; %%symmetriser leg
\draw (5,-2) -- (5,-2.5) ; %%symmetriser leg
\draw (6,-2) -- (6,-2.5) ; %%symmetriser leg
%\draw (4,-2) -- (3,-2.5) ; %%symmetriser leg
\draw (1.79,-2.5) -- (4.1,-2.5) -- (4.1,-3.20992) -- (1.79,-3.20992) -- (1.79,-2.5);   %% the symmetrizer 
\node at (3,-2.79) {$a$} ; 
%     \foreach \x in {#4,#5,#6,#7,#8}
%          \draw (\x,0) -- (\x,-1);  %% the vert lines 
%     \foreach \x % [evaluate \x] 
%          in {1,2,...,#4}
%     \draw (\x,0) -- (\x,-1);       %% the vert lines before cup
%     \foreach \x in {#7,...,#8}
%        \draw (\x,0) -- (\x,-1);    %% the vert lines after cup
%%%     \draw (2.2385,0) -- (2.2385,-1);
     \draw (#2,0) .. controls (#2,-.5) and (#3,-.5) .. (#3,0) ;
 \end{tikzpicture}
 }

